\begin{frame}{softmax를 다시 볼 이유: 차이를 증폭}
  Week 2.3(로지스틱 회귀의 출력을 확률로), Week 3(합이 1인
  분포)에서 만난 softmax가, 어텐션 맥락에서 더 예리하게 —
  \kq{\textbf{원점수 간 ``차이''를 지수적으로 증폭}한다.}
  \vskip0.4em
  \textbf{2개 단어의 가장 단순한 경우} ($s_1=0$, $s_2=g$):
  \[p_2 = \frac{e^{g}}{1 + e^{g}} = \frac{1}{1 + e^{-g}}
      = \sigma(g)\]
  \textbf{어텐션 원점수 차이에 대한 softmax는 정확히 시그모이드
  $\sigma(g)$} — 로지스틱 회귀의 결정 함수가 ``원점수
  \emph{차이}''를 입력으로 받는 모습. 두 원점수 사이의 ``차이''
  크기가 가중치의 극단적 정도를 사실상 \textbf{단독으로}
  결정한다:
  \small
  \begin{tabular}{@{}lll@{}}
    \toprule
    원점수 차이 $g$ & 가중치 $(p_1, p_2)$ & 해석 \\
    \midrule
    0.5 & (0.378, 0.622) & 아직 균형 \\
    2 & (0.119, 0.881) & 한쪽 우세 \\
    8 & (0.0003, 0.9997) & 사실상 하드 선택 \\
    \bottomrule
  \end{tabular}
  \vskip0.4em
  차이가 0.5에서 8로 \textbf{16배}가 되자 ``62:38의 균형''이
  ``99.97:0.03의 독주''로. 원점수가 큰 차원의 내적(차원에 비례해
  커짐)을 그대로 넣으면 \textbf{거의 one-hot} — ``관련 단어
  \emph{하나}만 고르는'' 하드 경로가 되어 (1) \textbf{선택된
  단어의 가중치 미분이 0}(학습 신호 끊김), (2) \textbf{원점수
  미세 변화에 가중치가 급격히 뒤집혀} 불안정.
  $\sqrt{d_k}$ 스케일링은 이 \textbf{``softmax의 극단화''를 막는
  안전장치}.
\end{frame}
